\chapter{Appendices}
\addcontentsline{toc}{chapter}{Appendices}

\section*{Appendix A: The Five Perception Modes in Operation}
\addcontentsline{toc}{section}{Appendix A: The Five Perception Modes in Operation}

The screenshots below show Echora's five modes running on real inputs: navigation with proximity-zoned obstacles, text reading, face recognition, banknote identification, and hand-guidance interaction.

\begin{figure}[H]
\centering
\begin{minipage}[t]{0.31\textwidth}\centering
  \includegraphics[width=\linewidth,height=3.4cm,keepaspectratio]{mode_navigation.jpg}\\[2pt]{\footnotesize (a) Navigation}
\end{minipage}\hfill
\begin{minipage}[t]{0.31\textwidth}\centering
  \includegraphics[width=\linewidth,height=3.4cm,keepaspectratio]{mode_ocr.jpg}\\[2pt]{\footnotesize (b) Text reading (OCR)}
\end{minipage}\hfill
\begin{minipage}[t]{0.31\textwidth}\centering
  \includegraphics[width=\linewidth,height=3.4cm,keepaspectratio]{mode_face.jpg}\\[2pt]{\footnotesize (c) Face recognition}
\end{minipage}
\\[6pt]
\begin{minipage}[t]{0.31\textwidth}\centering
  \includegraphics[width=\linewidth,height=3.4cm,keepaspectratio]{mode_banknote.jpg}\\[2pt]{\footnotesize (d) Banknote identification}
\end{minipage}\hfill
\begin{minipage}[t]{0.31\textwidth}\centering
  \includegraphics[width=\linewidth,height=3.4cm,keepaspectratio]{mode_interaction.jpg}\\[2pt]{\footnotesize (e) Hand guidance}
\end{minipage}\hfill
\begin{minipage}[t]{0.31\textwidth}\centering
  ~
\end{minipage}
\caption{Echora's five perception modes running on real inputs.}
\label{fig:modes}
\source{Teamwork}
\end{figure}

\section*{Appendix B: Hardware Build Highlights}
\addcontentsline{toc}{section}{Appendix B: Hardware Build Highlights}

The prototype was built by hand. Figure~\ref{fig:build} shows the finished 3D-printed headset and the custom haptic module that drives the wrist array of vibration motors.

\begin{figure}[H]
\centering
\begin{minipage}[b]{0.48\textwidth}\centering
  \includegraphics[width=\linewidth,height=5.2cm,keepaspectratio]{build_glasses.jpeg}\\[2pt]{\footnotesize (a) Finished headset}
\end{minipage}\hfill
\begin{minipage}[b]{0.48\textwidth}\centering
  \includegraphics[width=\linewidth,height=5.2cm,keepaspectratio]{build_haptic.png}\\[2pt]{\footnotesize (b) Custom haptic module}
\end{minipage}
\caption{Highlights of the physical prototype build.}
\label{fig:build}
\source{Teamwork}
\end{figure}

\section*{Appendix C: Banknote Model Precision--Recall Curve}
\addcontentsline{toc}{section}{Appendix C: Banknote Model Precision--Recall Curve}

\begin{figure}[H]
\centering
\includegraphics[width=0.72\textwidth,keepaspectratio]{banknote_pr.png}
\caption{Precision--recall curve for the custom Egyptian banknote detector across all eleven denomination classes.}
\label{fig:bank_pr}
\source{Teamwork}
\end{figure}

\section*{Appendix D: Record of Generative AI Use}
\addcontentsline{toc}{section}{Appendix D: Record of Generative AI Use}

This appendix records how generative AI was used in the preparation of this
report, as required by the module guide. The declaration in the preliminary
pages should be read together with this record.

\textbf{Tool used.} A general-purpose generative AI writing assistant, accessed
through its standard web interface.

\textbf{What it was used for.} The tool was used only after we had written a
complete first draft of each chapter ourselves. Its use was limited to language
and presentation:

\begin{itemize}
  \item checking spelling and correcting typing errors;
  \item correcting grammar, punctuation, and sentence structure;
  \item improving the clarity and flow of paragraphs we had already written, so
        that the writing reads in a more formal and professional academic style;
  \item checking that the report followed the required format and structure,
        including consistent headings, captions, and the ordering of sections.
\end{itemize}

\textbf{Typical way it was used.} We pasted a paragraph or section of our own
draft into the tool and asked it to correct the grammar and improve the clarity
without changing the meaning. We then compared its version with our own,
accepted only the changes we agreed with, and rewrote anything that did not
match what we meant. The final wording in this report was chosen by us.

\textbf{What it was not used for.} The tool was not used to generate the ideas,
technical content, system design, code, experimental results, analysis, or
conclusions presented in this report. It was not used to find, suggest, or
produce any reference. Every source in the reference list was located, read, and
checked by the team, and every result reported was produced by running our own
evaluation scripts on our own trained models.

\textbf{Verification.} We are aware that AI-generated text can contain errors,
invented references, or bias. For this reason every sentence that the tool
helped to reword was checked by us against the original draft and against the
underlying technical work, to confirm that the meaning was unchanged and
factually correct.
